\section{Defining $\cj$}
\label{sec:cajal}

% Programs have behavior
% Compiler preserves program behavior
% Programs compile to k-linear maps

\subsection{Syntax and Evaluation}

$\cj$\! is a linearly typed programming language. Fig. \ref{fig:syntax}
details its syntax, similar to a lambda calculus with sum and product types \citep{walker2005substructural}. 
The main differences are the absence of higher-order functions,\footnote{Adding higher-order functions is possible, but not
not necessary for the applications shown in this paper.} 
and the inclusion of dictionaries and 
nondeterministic choice.

% Syntax
\begin{figure}[ht]
  \begin{mdframed}[topline=false,innertopmargin=-1.19ex,innerleftmargin=-0.1ex,innerrightmargin=0ex,linewidth=.8pt,skipbelow=2ex]
    \rulediv{18pt}{\large \t{Syntax}}
    \vspace{.1em}
    \[
    \begin{tabular}{c c}
        \scalebox{.8}{
        $\begin{aligned}
           \nb{e} \,\,\,\coloneqq 
                \ & \ \ \ \nb{x} \in \nb{\t{Var}} &&\t{variable}\\
                &\mid \unit &&\t{unit}\\
                &\mid \nb{(e_1,e_2)} &&\t{tuple}\\
                &\mid \nb{\inj{i}{e}} &&\t{injection}\\
                &\mid \nb{\dict{e_\t{k}}{e_\t{v}}} &&\t{dictionary}\\
                &\mid \nb{e_1;e_2} &&\t{sequencing}\\
                &\mid \nb{\proj{i}{e}} &&\t{projection}\\
                &\mid \nb{\case{e}{x_1}{e_1}{x_2}{e_2}} &&\t{case analysis}\\
                &\mid \nb{e_1 \sqcap e_2} &&\t{nondeterminism}\\
                &\mid \nb{e_1(e_2)_\mathcal{R}} &&\t{lookup}\\
                &\mid \nb{\letb{x}{e_1}{e_2}} &&\t{let binding}\\[4mm]
        \end{aligned}$}
    &
        $\begin{array}{c}
        \scalebox{.8}{%                                         
        $\begin{aligned}                                         
            \nb{\tau} \,\,\,\coloneqq
                &\ \ \ \nb{\Unit} &&\t{unit}\\
                &\mid \nb{\ty_1 \x \ty_2} &&\t{product}\\
                &\mid \nb{\ty_1 \+ \ty_2} &&\t{sum}\\
                &\mid \nb{\ty_1 \rightharpoonup \ty_2} &&\t{dictionary}\\
        \end{aligned}$}                                          
        \\[13mm]
        \scalebox{.8}{%     
        $\begin{aligned}
            \nb{\D} \,\,\,\coloneqq
            &\ \ \ \nb{\emp} &&\t{empty}\\
            &\mid \nb{\D,\bind{x}{\ty}} &&\t{binder}\\
        \end{aligned}$}           
        \end{array}$

    \\[15mm]
    \end{tabular}
    \]
  \end{mdframed}

  \caption{Syntax of $\cj$expressions $\nb{e}$, types $\nb{\ty}$, and contexts $\nb{\D}$.}
  \label{fig:syntax}
\end{figure}

\paragraph{Atom} The simplest expression in $\cj$\! is $\unit$, which introduces the unit
type $\nb{\Unit}$. To use $\unit$, we use sequencing. For example, $\unit;\unit$ sequences two $\unit$ together, and just returns the
second. More generally, for an expression like $\nb{e_1;e_2}$, sequencing evaluates $\nb{e_1}$ and $\nb{e_2}$, then
returns the value of $\nb{e_2}$. There is a formal specification of this behavior in Fig. \ref{fig:evaluating}, which
defines how expressions evaluate.
\footnote{It is an inductively defined relation \citep{pierce2002types}. These are essentially implications.
If the statement above the line holds, you may conclude the statement below the line.}



% Evaluating
\begin{figure}[h]
  \begin{mdframed}[topline=false,innertopmargin=-1.39ex,innerleftmargin=-0.1ex,innerrightmargin=0ex,linewidth=.8pt,skipbelow=2ex]
    \rulediv{28pt}{\large \t{Evaluating}}
    \[  
      % \SetCell[c=2]{c}
    \begin{tabular}{c}
        \scalebox{.86}{
        $\step {e} {v} \iff \t{expression }\nb{e} \t{ evaluates to }\nb{v}$}
    \\[5mm]
        \scalebox{.8}{
        \def\extraVskip{3pt}
        \def\defaultHypSeparation{\hskip .05in}
        \bottomAlignProof
        \AxiomC{$\step{e}{v}$}
        \UnaryInfC{$\step{\inj{i}{e}}{\inj{i}{v}}$}
        \DisplayProof}
        
        \quad

        \scalebox{.8}{
        \def\extraVskip{3pt}
        \def\defaultHypSeparation{\hskip .05in}
        \bottomAlignProof
        \AxiomC{$\step{e_\t{k}}{v_\t{k}}$}
        \AxiomC{$\step{e_\t{v}}{v_\t{v}}$}
        \BinaryInfC{$\step{\dict{e_\t{k}}{e_\t{v}}}{\dict{v_\t{k}}{v_\t{v}}}$}
        \DisplayProof}

        \quad

        \scalebox{.8}{
        \def\extraVskip{3pt}
        \def\defaultHypSeparation{\hskip .05in}
        \bottomAlignProof
        \AxiomC{$\step{e}{(e_1,e_2)}$}
        \AxiomC{$\step{e_i}{v_i}$}
        \BinaryInfC{$\step{\proj{i}{e}}{v_i}$}
        \DisplayProof}

        \quad

        \scalebox{.8}{
        \def\extraVskip{3pt}
        \def\defaultHypSeparation{\hskip .05in}
        \bottomAlignProof
        \AxiomC{$\ess{i} \bstep \vsrc{i}$}
        \UnaryInfC{$\ess1 \sqcap \ess2 \bstep \vsrc{i}$} 
        \DisplayProof}

    \\[5mm]

        \scalebox{.8}{
        \def\extraVskip{3pt}
        \def\defaultHypSeparation{\hskip .05in}
        \bottomAlignProof
        \AxiomC{$\step{e_1}{\unit}$}
        \AxiomC{$\step{e_2}{v_2}$}
        \BinaryInfC{$\step{e_1;e_2}{v_2}$}
        \DisplayProof}

        \quad

        \scalebox{.8}{
        \def\extraVskip{3pt}
        \def\defaultHypSeparation{\hskip .05in}
        \bottomAlignProof
        \AxiomC{$\step {e_1} {v_1}$}
        \AxiomC{$\step {\{x \map v_1\}(e_2)} {v_2}$}
        \BinaryInfC{$\step {\letb{x}{e_1}{e_2}} {v_2}$}
        \DisplayProof}

        \quad

        \scalebox{.8}{
        \def\extraVskip{3pt}
        \def\defaultHypSeparation{\hskip .05in}
        \bottomAlignProof
        \AxiomC{$\step{e}{\inj{i}{v}}$}
        \AxiomC{$\step{\{x_i \map v\}(e_i)}{v_i}$}
        \BinaryInfC{$\step{\case{e}{x_1}{e_1}{x_2}{e_2}}{v_i}$}
        \DisplayProof}

    \\[5mm]


        \scalebox{.8}{
        \def\extraVskip{3pt}
        \def\defaultHypSeparation{\hskip .05in}
        \bottomAlignProof
        \AxiomC{$\step{e_\t{kv}}{\dict{v_{\t{k}}}{v_\t{v}}}$}
        \AxiomC{$\step{e_\t{q}}{v_\t{q}}$}
        \AxiomC{$\nb{\mathcal{V}} = \{\nb{(v_\t{v})_i} \mid (\nb{v_\t{q}},\nb{(v_\t{k})_i}) \in \nb{\mathcal{R}}\}$}
        \AxiomC{$\text{\Large$\sqcap$}_{\nb{\mathcal{V}}}\,\nb{v}\bstep \vsrc{\t{r}}$}
        \QuaternaryInfC{$\step{e_\t{kv}(e_\t{q})_\mathcal{R}}{v_\t{r}}$}
        \DisplayProof}

    \\[3mm]

    \end{tabular}
    \]
  \end{mdframed}
  \caption{Evaluating $\cj$expressions}
  \label{fig:evaluating}
\end{figure}



\paragraph{Character} 
% Without loss of generality we focus on an alphabet with 4 characters
% A | B | C | D
% Vowel | Embed1 | Embed2 | Embed3

On its own $\unit$ is not very exciting. However, $\unit$ becomes powerful when
paired with other expressions. With injections, $\unit$ can encode characters in an alphabet, and are how we introduce values
of sum type. For example, the following injections introduce values of type 
$\nb{\Unit \+ \Unit \+ \Unit \+ \Unit}$, the type of a four letter alphabet $\nb{\t{Char}}$.
For convenience, we use $\nb{n}$-ary forms of $\cj$\! expressions, which are encodable using their
binary form.

\[
\nb{\t{a}} \triangleq\inj{0}{\unit} \quad\quad
\nb{\t{b}} \triangleq\inj{1}{\unit} \quad\quad
\nb{\t{c}} \triangleq\inj{2}{\unit} \quad\quad
\nb{\t{d}} \triangleq\inj{3}{\unit}
\]

To use expressions of sum type we pattern match, specifying what to do when we match against
each character. For example, the following pattern match inspects $\nb{x}$, returning $\nb{\t{a}}$ if the
it is a vowel, and otherwise returning $\nb{\t{d}}$.

\[
\nb{\t{isVowel}} \!\quad\triangleq\quad\! \nb{\t{case}\,x\,\t{of}\,
  \begin{cases}
    \inj{0}{x_0}\Rightarrow x;\inj{0}{x_0}\\
    \inj{1}{x_1}\Rightarrow x;\inj{3}{x_0}\\
    \inj{2}{x_2}\Rightarrow x;\inj{3}{x_0}\\
    \inj{3}{x_3}\Rightarrow x;\inj{3}{x_0}\\
  \end{cases}}
  \simeq\quad
  \nb{\t{case}\,x\,\t{of}\,
  \begin{cases}
    \t{a}\Rightarrow \t{a}\\
    \t{b}\Rightarrow \t{d}\\
    \t{c}\Rightarrow \t{d}\\
    \t{d}\Rightarrow \t{d}
  \end{cases}}
\]

% isVowel: case analysis
% string: product intro
% string reverse: product elim

% Products allow us to encode n-ary sequences ok cool
% Dictionaries 
% -alternative to select/aggregate, more general by allowing aggregation on discrete
% types.

\paragraph{Sequence} Tuples introduce expressions of product type, which can encode sequences. For example, 
$\nb{(\t{a}, \t{b}, \t{c})}$ is a sequence of type $\nb{\t{Char} \x \t{Char} \x \t{Char}}$.
To use these product types, we use projections to witness their elements. For example,
$\nb{\pi_0(\t{a}, \t{b}, \t{c})}$ returns $\nb{\t{a}}$.

\paragraph{Dictionary} Dictionaries introduce expressions of dictionary type, which are
sequences of keys and values. For example, $\dict{(\t{a},\t{a},\t{b})}{(\t{a},\t{b},\t{c})}$
is a dictionary of type $\nb{\t{Char} \Map \t{Char}}$, comprised of three key-value pairs. To use this dictionary, we have a generalized
lookup operator, which matches a query against keys in the dictionary, returning a nondeterministic
choice between all matches. Since $\nb{\t{a}}$ maps both to $\nb{\t{a}}$ and $\nb{\t{b}}$, querying
the dictionary with $\nb{\t{a}}$ introduces nondeterminism. Interestingly, dictionaries compile
to linear attention \citep{katharopoulos2020transformers}. They are an abstract model of attention, from a programming language theoretic
view.

In general, any expression in $\cj$\! with a valid type annotation evaluates
to some value. In other words, it has behavior. This is critical to prove, since 
what it will mean to correctly compile $\cj$\! programs is to preserve their behavior
as $\rd{k}$-linear maps.

\begin{bthm}{(Programs have behavior)}{}
\vspace{.25em}
\t{If }$\p{\D}{e}{\ty}$ \t{and} $\nb{\D\vdash\nv}$ \t{, then } $\exists \nb{v}, \step{\sigma(e)}{v}$  
\end{bthm}

\textit{Proof sketch.} The argument is by induction on $\p{\D}{e}{\ty}$, the type annotation
of programs. Appendix \ref{app:theory-proofs} contails a detailed proof. In short, the argument
considers every way to construct a program $\p{\D}{e}{\ty}$, and in each case shows how substituting its free
variables with values $\nb{\sigma}(\nb{e})$ and then evaluating the resulting program will always yield a value $\nv(\nb{e}) \bstep \nb{v}$.

\subsection{Compiling}

To lift the closed program assumption, a compiler
must work on programs with free variables. These programs
specify only a part of the algorithm a system implements.
The eventual behavior of those free variables is not determined
by the programmer, but rather the target system---in a process known as 
\textit{linking} \citep{cardelli1997program}. In this work,
linking amounts to connecting the compiled $\rd{k}$-linear maps within
a neural network.

% Talk about the closed program assumption. 

% A correct compiler permits predictions about the behavior of compiled
% programs. By observing how a program behaves, we should be able to predict how
% a $\rd{k}$-linear map will behave on certain inputs.

% Compiling
\begin{figure}[ht]
  \begin{mdframed}[topline=false,innertopmargin=-1.39ex,innerleftmargin=-0.1ex,innerrightmargin=0ex,linewidth=.8pt,skipbelow=2ex]
    \rulediv{28pt}{\large \t{Compiling}}

    \[
    \scalebox{.86}{$\cm{-} :\nb{\t{Program}}\to \rd{k\t{-linear map}}$}\\[1em]
    \]
    \[
    \scalebox{.8}{
    $\begin{aligned}
    &\c{\p{\bind{x}{\ty}}{x}{\ty}}({\xvec}) = \xvec\\\\
    &\c{\p{\emp}{\unit}{\Unit}}(\rd{0}) = \rd{1}\\\\
    &\c{\p{\D}{\inj{1}{e}}{\ty_1 \+ \ty_2}}(\svec) = \rd{(}\cm{e}\rd{,}\,\rd{0}_{[\hspace{-.18em}[\nb{\ty_2}]\hspace{-.18em}]}\rd{)}\\\\
    &\c{\p{\D}{\inj{2}{e}}{\ty_1 \+ \ty_2}}(\svec) = \rd{(}\rd{0}_{[\hspace{-.18em}[\nb{\ty_2}]\hspace{-.18em}]}\rd{,}\,\cm{e}\rd{)}\\\\
    &\c{\p{\D}{(e_1,e_2)}{\ty_1 \x \ty_2}}(\svec) = \rd{(}\cm{e_1}\rd{,}\,\cm{e_2}\rd{)}\\\\
    &\c{\p{\D_\t{k} \o \D_\t{v}}{\dict{e_\t{k}}{e_\t{v}}}{\ty_\t{k} \Map \ty_\t{v}}}(\svec) = \xvec_\rd{\t{q}} \,\,\rd{\mapsto}\,\,\rsum_{\rd{i=1}}^{\nb{n}}\,\rd{\langle}\cm{\ess{\t{k}}}_\rd{i}\rd{,\,\xvec_\t{q}\rangle} \rt \cm{\ess{\t{v}}}_\rd{i}\\\\
    &\c{\p{\D_1 \o \D_2}{\letb{x}{e_1}{e_2}}{\ty_2}}(\svec) = \cm{e_2}(\cm{e_1})\\\\
    &\c{\p{\D_1 \o \D_2}{\case{e}{x_1}{e_1}{x_2}{e_2}}{\ty_2}}(\svec) = \cm{e_1}(\cm{e}_\rd{1}) \rp \cm{e_2}(\cm{e}_\rd{2})\\\\
    &\c{\p{\D}{\pi_ie}{\ty_i}}(\svec) = \cm{e}_\rd{i}\\\\
    &\c{\p{\D_\t{kv}\o\D_\t{q}}{e_\t{kv}(e_\t{q})_\mathcal{R}}{\ty_\t{v}}}(\svec) = \cm{e_\t{kv}}(\cm{\mathcal{R}}(\cm{e_\t{q}}))\\\\
    &\c{\p{\D}{e_1 \sqcap e_2}{\ty}}(\svec) = \cm{e_1}\rp\cm{e_2}\\\\
    &\c{\p{\D}{e_1;e_2}{\ty}}(\svec) = \cm{e_1}\rt\cm{e_2}\\\\[1em]
    \end{aligned}$}
    \]

  \end{mdframed}
  \caption{Compiling $\cj$ programs}
  \label{fig:compiling}
\end{figure}

% State in plain terms what the compiler is.
The compiler in Fig. \ref{fig:compiling} is a recursive function over programs.
In general, it compiles a program with $\rd{k}$
free variables into a $\rd{k}$-linear map. The domain and range of this
map is determined by the type annotations. For example, a program 
$\p{\bind{x}{\t{Char}}, \bind{x}{\t{Char}}}{e}{\t{Char}}$ compiles to a 
$\rd{2}$-linear map $\rd{f}: \rd{\R^4} \x \rd{\R^4} \to \rd{\R^4}$, with
$\cm{\t{Char}}=\rd{\R^4}$. To illustrate how the compiler works, we will compile a few examples.
There are additional examples across Appendices \ref{app:theory-examples} and 
\ref{app:theory-proofs}.

\paragraph{Atom} Because there is no way for $\unit$ to depend on a variable, it
compiles to a constant $\rd{0}$-linear map, always producing $\rd{1}$. It is the simplest
linear map the compiler produces. However, sequencing can depend on free variables,
and help illustrate the simplest $\rd{2}$-linear map.

\[
\c{\p{\bind{x}{\Unit},\bind{y}{\Unit}}{x;y}{\Unit}}(\rd{x},\rd{y})=\c{\p{\bind{x}{\Unit}}{x}{\Unit}}(\rd{x})\rt \c{\p{\bind{y}{\Unit}}{y}{\Unit}}(\rd{y})=\rd{x}\rd{y}
\]

\paragraph{Character} The encoding of characters from before compile to one-hot encodings.
Since they lack free variables, they compile to constant $\rd{0}$-linear maps.

\[
\c{\p{\emp}{\t{a}}{\t{Char}}}(\rd{0}) = 
\c{\p{\emp}{\inj{0}{\unit}}{\t{1}\+\t{1}\+\t{1}\+\t{1}}}(\rd{0}) =
\rd{(1,0_{\bl{\c{\nb{\t{1}\+\t{1}\+\t{1}}}}})} =
\rd{\begin{bmatrix}1 \!&\! 0 \!&\! 0 \!&\! 0\end{bmatrix}^{\t{T}}}
\]

With pattern matching, we can compile computation which conditions on character
identity. In general, pattern matching across $\nb{n}$ cases compiles
to a linear mixture of $\nb{n}$ experts, gated by the expression we are matching
against \cite{jacobs1991adaptive}. For example, $\nb{\t{isVowel}}$ compiles to a $\rd{1}$-linear map,
returning a linear mixture of $\nb{4}$ experts.\footnote{When clear from context, we omit the full type annotation when compiling.}
% In each paragraph try to anchor the explanation with an intuition from
% the machine learning community.

\[
\c{\nb{\t{isVowel}}}
\Bigl(\rd{\begin{bmatrix}x_0\\x_1\\x_2\\x_3\end{bmatrix}}\Bigr) 
=
\cm{x;\inj{0}{\unit}}(\rd{x_0}) \rp
\rd{\sum_{i\bl{=}1}^{3}}\; 
\c{\nb{x;\inj{3}{\unit}}}(\rd{x_i})
=
\rd{x_0 \cdot \rd{\begin{bmatrix}1\\0\\0\\0\end{bmatrix}}\;+\;}
\;\rd{\sum_{i\bl{=}1}^{3}}\;
\rd{x_i \cdot \rd{\begin{bmatrix}0\\0\\0\\1\end{bmatrix}}}
\]


\paragraph{Sequence} Compiling a sequence is just concatenating the compiled entries of the sequence. Projection
amounts to projecting out the $\rd{i}$-th piece of that concatenated vector.
\[
\c{\nb{\pi_1\t{(a,b)}}:\nb{\t{Char}}} = \c{\nb{\t{(a,b)}}:\nb{\t{Char} \x \t{Char}}}_\rd{1} = 
\rd{\begin{bmatrix}1 \!&\! 0 \!&\! 0 \!&\! 0 \!&\! 0 \!&\! 1 \!&\! 0 \!&\! 0\end{bmatrix}^{\t{T}}_1}
= \rd{\begin{bmatrix}0 \!&\! 1 \!&\! 0 \!&\! 0\end{bmatrix}^{\t{T}}}
\]

\paragraph{Dictionary} We model attention using dictionaries. These dictionaries compile to linear attention with
 an identity kernel \citep{katharopoulos2020transformers}, which is essentially a sum of outer products between keys and values. 
 Lookup amounts to the application of a compiled query against the linear attention map.

\vspace{-1em}
\[
\c{\dict{(\t{a},\t{a},\t{b})}{(\t{a},\t{b},\t{c})}} = 
\cm{\t{a}} \,\rd{\otimes}\, \cm{\t{a}} \rp 
\cm{\t{b}} \,\rd{\otimes}\, \cm{\t{a}} \rp 
\cm{\t{c}} \,\rd{\otimes}\, \cm{\t{b}} \;\simeq
  \rd{\begin{blockarray}{ccccc}
   & \nb{\t{a}} & \nb{\t{b}} & \nb{\t{c}} & \nb{\t{d}} \\                                                   
  \begin{block}{c[cccc]}
    \nb{\t{a}} & 0 & 0 & 0 & 0 \rule{0pt}{2.2ex}\\                                                      
    \nb{\t{b}} & 0 & 0 & 0 & 0 \\                                                      
    \nb{\t{c}} & 1 & 1 & 0 & 0 \\                                                      
    \nb{\t{d}} & 1 & 0 & 0 & 0 \rule[-1ex]{0pt}{0pt}\\                                                      
  \end{block}     
  \end{blockarray}}
\]

We know this compiled dictionary is \textit{correct} because it mirrors the behavior
of the original dictionary. A lookup on $\nb{\t{b}}$ returns $\nb{\t{c}}$, and
looking up $\cm{\t{b}}$ in the compiled dictionary returns $\c{\nb{\t{c}}}$. 
The key theoretical result in this paper proves that this sense of correctness holds for all programs in
$\cj\!$. 



\begin{bthm}{(Compiler preserves program behavior)}{}
\vspace{.25em}
\t{If }$\p{\D}{e}{\ty}$ \t{and} $\nb{\D\vdash\nv}$, \t{and} $\nb{V}=\{(\nb{D},\vs) \mid \nv(\es) \bstep_\nb{D} \vs\}$, \t{ then } $\cm{e}(\cm{\nv})=\rd{\sum_{\bl{(\nb{D}\hspace{.05em},\hspace{.05em}\vs)\in\nb{V}}}}\,\cm{v}$
\end{bthm}

\textit{Proof sketch.} The argument is again by induction on $\p{\D}{e}{\ty}$, the type annotation
of programs. Appendix \ref{app:theory-proofs} contails a detailed proof. In short, the argument
considers every way to construct a program $\p{\D}{e}{\ty}$, and in each case shows that
the compiled program $\cm{e}$ mirrors the behavior $\nb{V}$ of the original program. Because programs
have nondeterministic behavior, $\nb{V}$ captures the possible nondeterministic values. Beyond this, most cases
rely on the following lemma.

\begin{blem}{(Programs compile to $\rd{k}$-linear maps)}{}
  \vspace{.25em}
\t{If }$\p{\D,\bind{x}{\ty_1}}{e}{\ty_2}$ \t{ then } $\cm{e}(\svec,\a1\rt\xv1\rp\a2\rt\xv2)=\a1\rt\cm{e}(\svec,\xv1)\rp \a2\rt\cm{e}(\svec,\xv2)$
\end{blem}
\vspace{.4em}

This lemma formally establishes the correspondence between programs and $\rd{k}$-linear maps
we have discussed throughout this paper. It is also the most important technical idea
in this paper---it further shows the promise of linear types for compiling discrete
algorithms into various kinds of neural networks \citep{velez2026compiling,velez2026compiling2}. 